\section{硅谷御三家的路径依赖}

本期对 Anthropic、OpenAI、Gemini 的评价，核心不是八卦，而是路径依赖。每家公司上一阶段的成功，会塑造下一阶段的组织注意力。Anthropic 的 underdog 文化和 coding/agent 投入让它站上窗口；OpenAI 的 ChatGPT ToC 成功可能让它忽视 coding；Gemini 的生态和模型能力强，但战略上更像领先追随者。

\subsection{逐项对比}

\noindent\begin{longtable}{p{0.18\textwidth}p{0.25\textwidth}p{0.25\textwidth}p{0.22\textwidth}}
\toprule
公司 & 访谈中的优势 & 访谈中的风险 & 对 Coding 第二幕的启示 \\
\midrule
Anthropic & 重视细节、文化严格、Agent/Coding 窗口强 & 优势会被追赶，规模化仍有挑战 & 早投入工作流会形成窗口。 \\
OpenAI & 品牌、模型、产品和资本强 & ToC 成功可能成为战略惯性 & 过去成功会遮蔽新主线。 \\
Gemini & 生态、资源、模型能力强 & Coding 落后、声势不足 & 生态不能替代明确战略。 \\
Meta & 开源、资本、人才和挑战者位置 & 产品闭环不稳定 & 开源可能是后发追赶路径。 \\
xAI & 资源和关注度高 & 战略摇摆、团队流失 & 组织稳定性比明星叙事重要。 \\
\bottomrule
\end{longtable}

\begin{warningbox}{不要把公司判断当成永久结论}
本期是季度快照，不是长期排名。模型公司窗口期可能以月为单位变化。更稳定的分析方法，是看路径依赖：组织奖励什么、忽视什么、谁能更快拥抱新范式。
\end{warningbox}

\subsection{本章小结}

硅谷御三家的比较，本质是战略注意力比较。Coding 第二幕会奖励那些把模型、产品、数据和工程闭环放到同一条主线上的公司。

